\vspace{2.5pt plus 3pt minus 0.5pt}\noindent\hspace*{1.2em}\kbvnum{(2)}\hspace{0.55em}ಈ ವಾಕ್ಯಗಳಲ್ಲಿನ ಕ್ರಿಯಾಪದಗಳ ಪುರುಷ, ಲಿಂಗ, ವಚನ, ಕಾಲಗಳನ್ನು ತಿಳಿಸಿರಿ.\par

\vspace{2.5pt plus 3pt minus 0.5pt}\noindent\hspace*{1.2em}\kbvnum{{[}ಅ{]}}\hspace{0.55em}ಬ್ಯಾಂಕರನು ಹಣವನ್ನು ಸಾಲವಾಗಿ ಕೊಡುತ್ತಾನೆ.\par

\vspace{2.5pt plus 3pt minus 0.5pt}\noindent\hspace*{1.2em}\kbvnum{{[}ಬ{]}}\hspace{0.55em}ನಾವು ಈ ದಿನ ಊರಿನಿಂದ ಬಂದೆವು.\par

\vspace{2.5pt plus 3pt minus 0.5pt}\noindent\hspace*{1.2em}\kbvnum{{[}ಕ{]}}\hspace{0.55em}ಬೆಕ್ಕು ಇಲಿಯನ್ನು ತಿಂದಿತು.\par

\vspace{2.5pt plus 3pt minus 0.5pt}\noindent\hspace*{1.2em}\kbvnum{{[}ಡ{]}}\hspace{0.55em}ಗೌರಿಯು ಜಾತ್ರೆಗೆ ಹೋಗುವಳು.\par

\vspace{2.5pt plus 3pt minus 0.5pt}\noindent\hspace*{1.2em}\kbvnum{(3)}\hspace{0.55em}“ಜನರು ಭಕ್ತಿಯಿಂದ ಗಣಪತಿಯನ್ನು ಪೂಜೆಮಾಡುತ್ತಾರೆ.”\par

\indent — ಈ ವಾಕ್ಯದಲ್ಲಿ “ಪೂಜೆ” ಎಂಬ ನಾಮಪ್ರಕೃತಿಯನ್ನು ಕ್ರಿಯಾಪದವಾಗಿ ಬಳಸಿರಿ.\par

\vspace{2.5pt plus 3pt minus 0.5pt}\noindent\hspace*{1.2em}\kbvnum{(4)}\hspace{0.55em}ಈ ಕ್ರಿಯಾಪದಗಳನ್ನು ನಿಮ್ಮ ಸ್ವಂತ ವಾಕ್ಯಗಳಲ್ಲಿ ಬಳಸಿರಿ.\par

\indent ಮಾಡೆವು ಜಪಿಸಲಿ ಹೋದಾನು\par

\indent - ಬರೆಯಲಿ - ತಿಂದಿರಿ ಕೊಟ್ಟಿ\par

\vspace{2.5pt plus 3pt minus 0.5pt}\noindent\hspace*{1.2em}\kbvnum{(5)}\hspace{0.55em}“ಏಳು” ಧಾತುವಿನ ಭೂತಕಾಲ ಕ್ರಿಯಾಪದ ರೂಪಗಳನ್ನು ತಿಳಿಸಿರಿ.\par

\vspace{2.5pt plus 3pt minus 0.5pt}\noindent\hspace*{1.2em}\kbvnum{(6)}\hspace{0.55em}“ಗೋಪಾಲ ರೆಡ್ಡಿಯು ನಾಳೆಯ ದಿನ ಹಣವನ್ನು ಕೊಡುತ್ತಾನೆ.”—ಈ ವಾಕ್ಯದಲ್ಲಿರುವ ಕ್ರಿಯಾಪದದ ವಿಶೇಷವೇನು?\par

\vspace{2.5pt plus 3pt minus 0.5pt}\noindent\hspace*{1.2em}\kbvnum{(7)}\hspace{0.55em}ಕಾಲಸೂಚಕ ಪ್ರತ್ಯಯಗಳು ಯಾವುವು? ಉದಾಹರಣೆಗಳ ಸಹಿತ ತಿಳಿಸಿ.\par

\vspace{2.5pt plus 3pt minus 0.5pt}\noindent\hspace*{1.2em}\kbvnum{(8)}\hspace{0.55em}ವಿಧ್ಯರ್ಥಕ ಕ್ರಿಯಾಪದ ಎಂದರೇನು? ಒಂದು ಉದಾಹರಣೆ ಸಹಿತ ತಿಳಿಸಿ.\par

\vspace{2.5pt plus 3pt minus 0.5pt}\noindent\hspace*{1.2em}\kbvnum{(9)}\hspace{0.55em}ನಿಷೇಧಾರ್ಥಕ ಕ್ರಿಯಾಪದ ಎಂದರೇನು? ಉದಾಹರಣೆಗಳ ಸಹಿತ ತಿಳಿಸಿ.\par

\vspace{2.5pt plus 3pt minus 0.5pt}\noindent\hspace*{1.2em}\kbvnum{(10)}\hspace{0.55em}“ತರು” ಧಾತುವಿನ ಸಂಭಾವನಾರ್ಥಕ ಕ್ರಿಯಾಪದಗಳ ಏಕವಚನ ರೂಪಗಳನ್ನು ಬರೆದು ತಿಳಿಸಿ.\par

\vspace{2.5pt plus 3pt minus 0.5pt}\noindent\hspace*{1.2em}\kbvnum{(11)}\hspace{0.55em}“ಮತಿ - ಉಂಟು, ಗತಿ - ಇಲ್ಲ”-ಇಲ್ಲಿರುವ ಕ್ರಿಯಾಪದಗಳನ್ನು ಗುರ್ತಿಸಿ, ಅವುಗಳ ವೈಶಿಷ್ಟ್ಯ ನಿರೂಪಣೆಯನ್ನು ತಿಳಿಸಿರಿ.\par
